% =====================================================================
%  Chương 4 — Mục 4 của đề: một CNN cơ bản và ba mô hình phát triển
%  Mã nguồn: a05/models.py, a05/train.py, a05/experiment.py
%  Notebook: 03_cifar10.ipynb, 04_cifar100.ipynb, 05_diabetes.ipynb
% =====================================================================
\chapter{Một CNN cơ bản và ba mô hình phát triển}
\label{ch:mohinh}

Chương này trả lời mục 4 của đề: xây một CNN cơ bản và ba mô hình phát triển, rồi huấn luyện cả bốn trên ba tập
dữ liệu của Chương~\ref{ch:dulieu}. Cả bốn mô hình được viết từ đầu bằng PyTorch~\cite{paszke2019}, không dùng trọng số
huấn luyện sẵn.

\section{Nguyên tắc thiết kế: mỗi bước thêm đúng một cơ chế}

Nếu bốn mô hình khác nhau ở nhiều chỗ cùng lúc thì không thể biết chênh lệch kết quả đến từ đâu. Vì vậy mỗi mô hình
được dựng từ mô hình trước bằng cách thêm \emph{một} cơ chế ở Chương~\ref{ch:kientruc}:

\begin{itemize}
  \item \textbf{M0 BasicCNN}: khuôn LeNet/AlexNet, ba lần lặp Conv--ReLU--MaxPool rồi hai tầng nối đầy đủ. Không
        BatchNorm, không đường tắt.
  \item \textbf{M1 VGGNet}: ba tầng, mỗi tầng hai khối, mỗi khối là hai conv $3\times3$ có BatchNorm~\cite{ioffe2015};
        cuối mạng là Global Average Pooling thay cho các tầng nối đầy đủ lớn.
  \item \textbf{M2 ResNet}: đúng khung của M1, mỗi khối thêm đường tắt $y=\mathrm{ReLU}(F(x)+x)$. Khi số kênh đổi,
        đường tắt đi qua một conv $1\times1$.
  \item \textbf{M3 SE-ResNet}: đúng M2, thêm khối SE trên nhánh $F$ trước phép cộng.
\end{itemize}

Bước M0$\to$M1 gộp ba thay đổi (sâu hơn, BatchNorm, GAP), vì đó chính là bước chuyển từ AlexNet sang VGG trong lịch sử.
Hai bước sau thì mỗi bước chỉ khác đúng một cơ chế: M1 và M2 chỉ khác nhau ở đường tắt, M2 và M3 chỉ khác nhau ở khối SE.
Hình~\ref{fig:kientruc} đặt bốn mô hình cạnh nhau. Shape ghi bên phải mỗi khối được đọc từ một lượt forward thật trong
notebook 06, không phải tính tay.

\begin{figure}[H]
  \centering
  \hinh{f4_kien_truc}
  \caption{Bốn mô hình cho ảnh $32\times32\times3$; M1, M2, M3 dùng chung khung và chỉ khác nhau ở nội dung của khối}
  \label{fig:kientruc}
\end{figure}

\section{Mã nguồn bốn mô hình}

M0 là một lớp riêng. M1, M2, M3 dùng chung lớp \texttt{StagedCNN}, chỉ khác hai cờ \texttt{residual} và \texttt{se}
của lớp \texttt{Block}. Nhờ vậy ``chỉ khác một cơ chế'' được đảm bảo ngay trong mã nguồn: hai mô hình liền kề chạy qua
cùng một đoạn mã, chỉ khác một nhánh \texttt{if}.

\maNguon{models__BasicCNN}{M0: CNN cơ bản}
\maNguon{models__Block}{Khối hai conv $3\times3$, bật tắt đường tắt và SE bằng cờ}
\maNguon{models__StagedCNN}{Khung chung của M1, M2, M3}

Hàm \texttt{build\_model} dựng mô hình theo tên và theo số chiều. Tham số \texttt{dim} chọn giữa bộ
(\texttt{Conv2d}, \texttt{BatchNorm2d}, \texttt{MaxPool2d}) và bộ 1D tương ứng, nên cùng một định nghĩa phục vụ cả ảnh
lẫn dữ liệu bảng. Mọi tầng tích chập được khởi tạo He (\texttt{kaiming\_normal\_}), phù hợp với ReLU.

\maNguon{models__build_model}{Dựng mô hình theo tên, cho cả 2D và 1D}

\section{Dò shape tầng theo tầng}

Bảng~\ref{tab:traceM0} là đầu ra thật của từng tầng con khi một tensor $1\times3\times32\times32$ đi qua M0. Ba lần
MaxPool đưa $32\times32$ về $4\times4$; tầng Flatten nối $128\cdot4\cdot4$ giá trị vào tầng nối đầy đủ, và riêng tầng này
chiếm phần lớn số tham số của M0. Ở M1, M2, M3, Global Average Pooling nén mỗi kênh về một số, nên tầng nối đầy đủ cuối
chỉ còn $256\times K$ trọng số (Bảng~\ref{tab:traceM3}).

\begin{table}[H]
  \centering
  \begin{minipage}[t]{0.47\textwidth}
    \centering
    \caption{Dò shape của M0 (2D, $K=10$)}
    \label{tab:traceM0}
    \small\input{generated/tab_trace_basic}
  \end{minipage}\hfill
  \begin{minipage}[t]{0.47\textwidth}
    \centering
    \caption{Dò shape của M3 (2D, $K=10$)}
    \label{tab:traceM3}
    \small\input{generated/tab_trace_seresnet}
  \end{minipage}
\end{table}

Hình~\ref{fig:fmapm0} cho thấy tầng tích chập đầu của M0 (sau khi huấn luyện trên CIFAR-10) làm gì với một ảnh test: mỗi
kênh là ảnh vào sau khi qua một bộ lọc $3\times3\times3$. Có kênh làm nổi biên theo một hướng, có kênh phản ứng với
vùng sáng hoặc với một màu. Đây là đầu vào của các tầng sau.

\begin{figure}[H]
  \centering
  \hinh{f4_fmap_m0}
  \caption{Ảnh test lớp horse và 8 feature map đầu tiên sau tầng tích chập đầu của M0}
  \label{fig:fmapm0}
\end{figure}

\section{Bản 1D cho Diabetes}

Với Diabetes, \texttt{build\_model(name, dim=1, ...)} giữ nguyên bốn cơ chế, chỉ đổi mọi tầng 2D sang 1D, giảm độ rộng
xuống 16/32/64 kênh và cho ra một logit. Đầu vào là \texttt{[B, 1, \SL{data/diabetes/n_features_encoded}]}.
Bảng~\ref{tab:trace1d} là shape của M0 bản 1D; ba lần MaxPool làm chiều dài giảm từ \SL{data/diabetes/n_features_encoded}
xuống 1.

\begin{table}[H]
  \centering
  \caption{Dò shape của M0 bản 1D cho Diabetes}
  \label{tab:trace1d}
  \small\input{generated/tab_trace_basic1d}
\end{table}

\section{Huấn luyện}

\subsection{Cấu hình}

Trong cùng một tập, bốn mô hình dùng đúng một cấu hình (Bảng~\ref{tab:cauhinh}) và cùng hạt giống 42. Mỗi mô hình chạy
một lần. Checkpoint được giữ ở epoch có loss validation nhỏ nhất và được nạp lại trước khi đánh giá trên tập test.

\begin{table}[H]
  \centering
  \caption{Cấu hình huấn luyện}
  \label{tab:cauhinh}
  \small
  \begin{tabular}{lll}
    \toprule
    & CIFAR-10 và CIFAR-100 & Diabetes \\
    \midrule
    Thiết bị & GPU \SL{cifar10/device}, AMP fp16 & CPU, \SL{diabetes/config/threads} luồng \\
    Tối ưu & SGD nesterov, momentum \SL{cifar10/config/momentum} & AdamW \\
    Tốc độ học & OneCycle, đỉnh \SL{cifar10/config/max_lr} & cố định \SL{diabetes/config/lr} \\
    Weight decay & \SL{cifar10/config/weight_decay} & \SL{diabetes/config/weight_decay} \\
    Epoch, cỡ lô & \SL{cifar10/config/epochs}, \SL{cifar10/config/batch_size} &
                   tối đa \SL{diabetes/config/epochs}, \SL{diabetes/config/batch_size} \\
    Dừng sớm & không & sau \SL{diabetes/config/patience} epoch không giảm val loss \\
    Cắt gradient & $\lVert\nabla\rVert_2\le\SL{cifar10/config/grad_clip}$ & không \\
    Hàm mất mát & cross-entropy & BCE, \texttt{pos\_weight} = số âm / số dương \\
    Ngưỡng phân loại & argmax & chọn trên validation theo F1 \\
    \bottomrule
  \end{tabular}
\end{table}

Lịch OneCycle~\cite{smith2019} tăng tốc độ học từ nhỏ lên đỉnh trong 30\,\% số bước đầu, rồi giảm dần về gần 0.
Pha giảm cuối giúp mô hình hội tụ vào một điểm tốt trong \SL{cifar10/config/epochs} epoch mà không cần dò lịch giảm bằng tay.

\subsection{Vòng lặp huấn luyện}

Một epoch là một lượt qua dữ liệu: với mỗi lô, tính $L(f_\theta(x), y)$, lan truyền ngược, rồi cập nhật $\theta$. Hàm
\texttt{train\_one\_epoch} dùng chung cho cả đánh giá: khi \texttt{opt=None}, nó chỉ tính loss và accuracy. Với AMP, phép
tính chạy ở fp16, còn \texttt{GradScaler} nhân loss lên trước khi lan truyền ngược để gradient nhỏ không bị làm tròn về 0.

\maNguon{train__train_one_epoch}{Một epoch huấn luyện hoặc đánh giá}
\maNguon{train__fit}{Huấn luyện nhiều epoch, giữ checkpoint theo val loss}

Hai hàm dưới đây gói phần riêng của từng loại dữ liệu: bộ tối ưu, lịch tốc độ học, hàm mất mát và cách sinh lô.

\maNguon{train__train_image_model}{Huấn luyện một mô hình 2D trên CIFAR}
\maNguon{train__train_tabular_model}{Huấn luyện một mô hình 1D trên Diabetes}

\subsection{Sự cố phân kỳ của M0 và vai trò của BatchNorm}

Lượt chạy đầu tiên của M0 trên CIFAR-10 chưa cắt gradient và đã phân kỳ: loss trở thành NaN ở epoch
\SL{compare/noclip_diverged_epoch}, đúng lúc lịch OneCycle đưa tốc độ học lên đỉnh
\SL{cifar10/config/max_lr} (Hình~\ref{fig:clip}). Giải thích hợp lý nhất là M0 thiếu BatchNorm: BatchNorm chuẩn hoá đầu ra
mỗi tầng về trung bình 0 và phương sai 1, nên độ lớn của tín hiệu và của gradient không tăng dần qua các tầng, và vì thế
chịu được tốc độ học lớn hơn~\cite{ioffe2015}. Đây là suy luận từ cơ chế, không phải kết quả của một thí nghiệm tách riêng.

\begin{figure}[H]
  \centering
  \hinh{f4_clip}
  \caption{M0 trên CIFAR-10: không cắt gradient thì phân kỳ khi tốc độ học lên đỉnh}
  \label{fig:clip}
\end{figure}

Để bốn mô hình vẫn dùng chung một cấu hình, lượt chạy chính thức bật cắt chuẩn gradient cho \emph{cả bốn}. Bài không
chạy riêng M1, M2, M3 khi tắt cắt gradient, cũng không đo ngưỡng bị chạm bao nhiêu lần, nên không kết luận được việc cắt
có ảnh hưởng tới ba mô hình này hay không. Lượt phân kỳ được giữ lại trong
\texttt{cifar10\_basic\_noclip\_history.dat} làm bằng chứng.

\subsection{Một lượt thí nghiệm}

Notebook 03, 04, 05 đều chỉ gọi một hàm: huấn luyện đủ bốn mô hình, lưu checkpoint và \texttt{metadata.json}, ghi đường
học ra \texttt{.dat}, rồi đo chất lượng và chi phí trên tập test. Ba tập vì thế đi qua đúng cùng một quy trình.

\maNguon{experiment__run_image_experiment}{Một lượt thí nghiệm đầy đủ trên một tập ảnh}

\paragraph{Ghi chú về thời gian.} Máy dùng GPU laptop \SL{cifar10/device}. Trong suốt các lượt huấn luyện ảnh, GPU chạy
nóng và tự hạ xung nhịp vì nhiệt (\texttt{nvidia-smi} báo cờ \emph{SW thermal slowdown}). Số giây mỗi epoch ở
Chương~\ref{ch:sosanh} vì vậy chậm hơn nhiều so với lúc máy còn mát. Chúng vẫn so sánh được với nhau vì cả tám lượt chạy
trong cùng điều kiện, nhưng không nên đọc như thời gian chuẩn của loại GPU này.
